% Generated draft: requires mathematical and source review.
\documentclass[11pt]{article}
\usepackage[margin=25mm]{geometry}
\usepackage{fontspec}
\setmainfont{Latin Modern Roman}
\usepackage{amsmath,amssymb,mathtools,unicode-math}
\setmathfont{Latin Modern Math}
\usepackage{enumitem,needspace,xurl,hyperref}
\setlength{\parindent}{0pt}
\setlength{\parskip}{5pt}
\setlength{\emergencystretch}{3em}
\newcommand{\N}{\mathbb{N}}
\newcommand{\Z}{\mathbb{Z}}
\newcommand{\Q}{\mathbb{Q}}
\newcommand{\R}{\mathbb{R}}
\newcommand{\C}{\mathbb{C}}
\newcommand{\F}{\mathbb{F}}
\newcommand{\A}{\mathbb{A}}
\newcommand{\PP}{\mathbb{P}}
\newcommand{\E}{\mathbb{E}}
\newcommand{\eps}{\varepsilon}
\newcommand{\dd}{\,\mathrm{d}}
\begin{document}
\section*{UnsolvedMath ID: 30004575}
\textbf{Source problem code:} catalogue entry
\subsection{Definitions and mathematical statement}
For notation only, let \(n,d\) be positive integers and let \(B^H=(B^{H,1},\ldots,B^{H,d})\) denote a fractional Brownian motion with Hurst parameter \(H\in(0,1)\). Each standard component is a centered Gaussian process with covariance
\[
\E[B^{H,i}_tB^{H,i}_s]=\frac12\bigl(t^{2H}+s^{2H}-|t-s|^{2H}\bigr),\qquad s,t\geq0.
\]
The source does not specify the admissible Hurst parameters or the covariance between different components. A schematic equation expressing the requested extension is
\[
\dd X_t=b(X_t)\,\dd t+\sum_{i=1}^{d}V_i(X_t)\,\dd\mathbf B^{H,i}_t,\qquad X_0=x\in\R^n,
\]
where \(b\) denotes a drift vector field, \(V_1,\ldots,V_d\) denote noise vector fields, and \(\mathbf B^H\) denotes a suitable rough-path enhancement of the driver. This display is explanatory notation, not a specification of hypotheses supplied by the source. Its interpretation requires a choice of enhancement, coefficient assumptions, and a time interval on which the solution exists; none is specified in the record. The source also does not specify whether the intended drift may depend on time.

When the law of \(X_t\) has a density with respect to Lebesgue measure, write
\[
\mathbb P(X_t\in A)=\int_A p_t(x,y)\,\dd y
\]
for Borel sets \(A\subseteq\R^n\). Lebesgue density is the interpretation used here to explain the question; the source does not explicitly identify the state space or reference measure.

\textbf{Source question.} Can the sharp density-estimate results for fractional-Brownian rough differential equations be extended to equations with drift, despite the loss of fractional-Brownian scaling and the required sub-Riemannian accessibility conditions?

\textbf{Unresolved formulation.} The intended task is to extend the referenced sharp density estimates to a drift-containing equation, with appropriate accessibility hypotheses. The supplied record does not state the original estimates, the meaning of sharpness or locality, the coefficient class, the accessibility condition, the relevant geometric distance or control cost, or the quantifiers and parameter dependence of the proposed bounds. Consequently, it does not determine a unique precise inequality or fully quantified conjecture. In particular, no specific bracket-generating condition, positivity assertion, or upper or lower density bound is asserted in this draft.
\subsection{Short English statement}
The question asks whether sharp density estimates for fractional-Brownian rough differential equations extend to equations with drift, while accounting for the stated loss of fractional-Brownian scaling and the necessary sub-Riemannian accessibility conditions. The supplied record labels the problem partially solved, but leaves the target estimates and essential hypotheses unspecified. A precise theorem statement therefore requires further source clarification and human mathematical review.
\subsection{Sources}
\begin{itemize}
\item Supplied local UnsolvedMath registry record, permanent ID \(30004575\), title: \emph{Sharp Fractional Brownian Density Estimates with Drift}; exact source problem code: \texttt{catalogue entry}. Collection: \emph{Oberwolfach Reports Open Problems}. The record gives proposed year 2020 and does not identify an individual proposer.
\item Supplied report citation: \emph{New Directions in Rough Path Theory} (2021), Oberwolfach Reports, pp.\ 1955--2019, DOI \texttt{10.4171/owr/2020/40}. \url{https://doi.org/10.4171/owr/2020/40}. The background attributes the question to workshop open-problem discussions documented in Oberwolfach Reports; no problem-specific page or passage is supplied.
\item Supplied literature source: \emph{Density for rough equations with irregular drift}. \url{https://doi.org/10.1016/j.bulsci.2021.103011}.
\item Supplied literature source: Geng, Ouyang and Tindel, described in the record as \emph{hypoelliptic estimates}. \url{https://arxiv.org/abs/2008.01229}.
\item Supplied literature source: \emph{OWR 2020/40 report}. \url{https://doi.org/10.4171/owr/2020/40}.
\item Literature assessment dated exactly \texttt{2026-08-22}: the supplied curation labels the problem \texttt{partially\_solved} and reports density results for some rough differential equations with drift, including irregular or linear-growth drift classes, while describing the requested sharp local density estimates for the general drift case as remaining open. This is supplied dated curation derived from the OWR material, not an independently verified literature assessment. The record itself requests independent verification of time-sensitive status and cited claims.
\end{itemize}
\end{document}
